\documentclass[12pt,a4paper]{article}
\usepackage[utf8]{inputenc}
\usepackage[T1]{fontenc}
\usepackage[polish]{polski}
\usepackage{amsmath}
\usepackage{amsfonts}
\usepackage{amssymb}
\usepackage{booktabs}
\usepackage{geometry}
\usepackage{enumitem}
\usepackage{graphicx}
\usepackage{float}
\usepackage{listings}
\usepackage{xcolor}
\usepackage{tabularx}
\usepackage{array}
\usepackage{etoolbox}

% lstlisting nigdy nie łamany między stronami
\BeforeBeginEnvironment{lstlisting}{\begin{minipage}{\linewidth}}
\AfterEndEnvironment{lstlisting}{\end{minipage}}

\geometry{margin=2cm}

\definecolor{codebg}{rgb}{0.97,0.97,0.97}
\definecolor{codekey}{rgb}{0.0,0.3,0.7}
\definecolor{codecomment}{rgb}{0.2,0.5,0.2}
\definecolor{codestring}{rgb}{0.6,0.1,0.1}

\lstset{
    language=C++,
    basicstyle=\ttfamily\small,
    keywordstyle=\color{codekey}\bfseries,
    commentstyle=\color{codecomment}\itshape,
    stringstyle=\color{codestring},
    backgroundcolor=\color{codebg},
    frame=single,
    framesep=4pt,
    breaklines=true,
    showstringspaces=false,
    numbers=left,
    numberstyle=\tiny\color{gray},
    xleftmargin=2em,
    tabsize=4,
    columns=fullflexible,
    literate=
        {ą}{{\k{a}}}1 {Ą}{{\k{A}}}1
        {ę}{{\k{e}}}1 {Ę}{{\k{E}}}1
        {ó}{{\'o}}1   {Ó}{{\'O}}1
        {ś}{{\'s}}1   {Ś}{{\'S}}1
        {ż}{{\.z}}1   {Ż}{{\.Z}}1
        {ź}{{\'z}}1   {Ź}{{\'Z}}1
        {ć}{{\'c}}1   {Ć}{{\'C}}1
        {ń}{{\'n}}1   {Ń}{{\'N}}1
        {ł}{{\l{}}}1  {Ł}{{\L{}}}1,
}

\title{Systemy czasu rzeczywistego, Projekt 4}
\author{Piotr Gugnowski, wydział EiTI, nr indeksu 320690}
\date{Czerwiec 2026}

\pagestyle{empty}

\begin{document}

\maketitle

% ═══════════════════════════════════════════════════════════════════════════════
\section*{Zadanie 4.1 — System CAN: suma kontrolna}
% ═══════════════════════════════════════════════════════════════════════════════

\subsection*{Treść zadania}

Jak wiemy, determinizm jest bardzo ważnym, ale nie jedynym i~często nie
najważniejszym wymaganiem stawianym sieciom czasu rzeczywistego. W~praktyce
ogromną rolę odgrywa bezpieczeństwo przesyłanych informacji. Pod tym pojęciem
rozumiemy w~węższym sensie zgodność informacji nadawanej i~odbieranej w~sieci.
W~szerszym sensie bezpieczeństwo informacji przesyłanych w~sieci polega również
na ich odpowiednim zabezpieczeniu w~celu utrudnienia lub uniemożliwienia dostępu
do nich stronom trzecim, lub/i~zastosowaniu odpowiednich technik kryptograficznych.
W~tym zadaniu interesować nas będzie wyłącznie bezpieczeństwo w~zdefiniowanym
wyżej węższym sensie. Jak wiadomo, na skutek zakłóceń komunikacyjnych np.\
wskutek oddziaływania silnych zaburzeń elektromagnetycznych, przesyłana
informacja może ulec zniekształceniu. Dla celów detekcji błędów transmisji
stosowane są różne techniki. Jedną z~nich jest stosowana dość powszechnie
w~profesjonalnych sieciach komunikacyjnych redundantna informacja zwana sumą
kontrolną lub cykliczną sumą redundancyjną. Suma kontrolna jest tworzona według
określonego algorytmu na podstawie wysyłanych informacji przez nadajnik.
Suma ta dołączana jest do strumienia wysyłanej informacji. Odbiornik tworzy własną
sumę kontrolną według tego samego algorytmu co nadajnik, ale na podstawie
strumienia danych odbieranych. Różnica w~odebranej i~zrekonstruowanej w~odbiorniku
sumie kontrolnej wskazuje na wystąpienie błędu lub błędów komunikacyjnych
i~pozwala na podjęcie odpowiednich akcji. Należy jednak pamiętać, że:

\begin{itemize}
    \item istnieje skończone prawdopodobieństwo zgodności sum kontrolnych odebranej
          i~utworzonej w~odbiorniku mimo wystąpienia błędu lub błędów transmisji;
    \item prawdopodobieństwo to jest w~profesjonalnych systemach sieciowych znikome
          i~minimalizowane dzięki odpowiedniemu doborowi algorytmów wyznaczania sumy
          kontrolnej i~starannemu doborowi wielomianów generacyjnych.
\end{itemize}

W~przypadku sieci CAN stosowany jest następujący algorytm.

\subsubsection*{Algorytm wyznaczania CRC}

Niech \texttt{CRC\_RG(14:0)} oznacza rejestr przechowujący wartość sumy kontrolnej CRC,
natomiast \texttt{NXTBIT} oznacza kolejny bit transmitowanego ciągu bitów począwszy od
SOF i~skończywszy na ostatnim bicie pola DATA. Suma CRC obliczana jest następująco:

\begin{lstlisting}[language={},basicstyle=\ttfamily\small,numbers=none,frame=single,backgroundcolor=\color{codebg}]
CRC_RG = 0;
REPEAT
    CRCNXT = NXTBIT EXOR CRC_RG(14)
    CRC_RG(14:1) = CRC_RG(13:0);
    CRC_RG(0) = 0;
    IF CRCNXT THEN
        CRC_RG(14:0) = CRC_RG(14:0) EXOR (4599hex);
    ENDIF
UNTIL (zostanie osiągnięte CRC SEQUENCE lub wystąpi błąd transmisji)
\end{lstlisting}

Napisać aplikację wyznaczającą sumę kontrolną dla sieci CAN. Aplikacja ta ma być
optymalizowana ze względu na minimalizację czasu realizacji algorytmu wyznaczania CRC.

\textbf{Założenia:}
\begin{enumerate}[label=\alph*)]
    \item Aplikacja ma umożliwiać wprowadzenie z~klawiatury ciągu do 96~bitów.
    \item Aplikacja powinna wyznaczyć sumę kontrolną z~wprowadzonego ciągu
          bitowego i~wyświetlić ją w~postaci heksadecymalnej.
    \item Aplikacja powinna mieć możliwość deklarowania liczby powtórzeń obliczenia
          sumy kontrolnej od~1 do~$10^9$.
    \item Aplikacja powinna mieć możliwość pomiaru i~wyświetlania łącznego czasu
          realizacji zadeklarowanej liczby powtórzeń obliczenia sumy kontrolnej oraz
          średniego czasu jednokrotnej realizacji obliczenia CRC.
    \item Aplikacja powinna być dostarczona w~postaci kodu maszynowego wykonywalnego.
\end{enumerate}

\subsection*{Implementacja}

% TODO: uzupełnić po zaimplementowaniu

\subsection*{Wyniki i~pomiary czasu}

% TODO: uzupełnić po zaimplementowaniu

% ═══════════════════════════════════════════════════════════════════════════════
\section*{Zadanie 4.2 — System CAN: identyfikatory wiadomości}
% ═══════════════════════════════════════════════════════════════════════════════

\subsection*{Treść zadania}

Jak wiadomo, węzły CAN są generatorami wiadomości. Każda wiadomość ma swój
własny identyfikator. Zależnie od specyfikacji sieci CAN występują w~specyfikacji
standardowej identyfikatory 11-bitowe, a~w~specyfikacji rozszerzonej identyfikatory
29-bitowe. W~związku z~tym teoretycznie możliwe jest zaadresowanie maksymalnie
odpowiednio $2^{11}$ lub $2^{29}$ wiadomości.

Czy tak jest w~rzeczywistości? Proszę podać istniejące ograniczenia w~specyfikacji
CAN (o~ile występują).

\subsection*{Odpowiedź}

% TODO: uzupełnić — patrz Dodatek A, sekcja A.4

% ═══════════════════════════════════════════════════════════════════════════════
\section*{Zadanie 4.3 — System CAN: rzeczywisty czas trwania transmisji ramki}
% ═══════════════════════════════════════════════════════════════════════════════

\subsection*{Treść zadania}

Dokładna analiza czasu transmisji ramki CAN wskazuje, że przy transmisji tej samej
liczby bajtów informacji w~polu danych czas trwania transmisji jest różny. Czas ten
zależy od kontekstu. Przyczyna leży w~sposobie kodowania informacji, a~w~szczególności
w~zastosowanej technice synchronizacji transferów sieciowych wspieranych przez
dodatkowe bity synchronizujące (technika \textit{bit stuffing}).

\textbf{Założenia:}
\begin{enumerate}[label=\alph*)]
    \item Prędkość transmisji w~sieci CAN: 1\,Mb/s.
    \item Przesyłamy w~polu danych 8~bajtów informacji.
    \item Identyfikator wiadomości podlega swobodnemu wyborowi z~ew.\ ograniczeniami
          (por.\ zadanie~4.2).
    \item Stosujemy ramkę o~identyfikatorze 29-bitowym.
    \item Współczynnik efektywności transmisji zdefiniowany jako stosunek liczby bitów
          wysłanej informacji (64~bity) do całkowitej liczby bitów ramki (łącznie
          z~bitami przerwy między ramkami i~bitami dodatkowymi dodanymi przez operację
          \textit{bit stuffing}).
\end{enumerate}

\begin{enumerate}[label=\alph*)]
    \item Wyznaczyć najkrótszy czas transmisji całej ramki CAN. Podać przykład
          zawartości pola danych, dla której ma to miejsce. Podać wartość współczynnika
          efektywności transmisji.
    \item Wyznaczyć najdłuższy czas transmisji całej ramki CAN. Podać przykład
          zawartości pola danych, dla której ma to miejsce. Podać wartość współczynnika
          efektywności transmisji.
\end{enumerate}

\subsection*{a) Najkrótszy czas transmisji}

% TODO: uzupełnić — patrz Dodatek B, sekcja B.5

\subsection*{b) Najdłuższy czas transmisji}

% TODO: uzupełnić — patrz Dodatek B, sekcja B.5


% ═══════════════════════════════════════════════════════════════════════════════
%  APPENDIX
% ═══════════════════════════════════════════════════════════════════════════════
\appendix

% ━━━━━━━━━━━━━━━━━━━━━━━━━━━━━━━━━━━━━━━━━━━━━━━━━━━━━━━━━━━━━━━━━━━━━━━━━━━━
\section*{Dodatek A — Tło teoretyczne}
% ━━━━━━━━━━━━━━━━━━━━━━━━━━━━━━━━━━━━━━━━━━━━━━━━━━━━━━━━━━━━━━━━━━━━━━━━━━━━

% ─────────────────────────────────────────────────────────────────────────────
\subsection*{A.0 Wyjaśnienia pojęć z~sekcji A.1}

\textbf{Sygnałowanie różnicowe — skąd nazwa?}

Zwykłe kable cyfrowe (np.\ USB w~trybie single-ended) przesyłają napięcie
mierzone względem masy: \texttt{0\,V} = bit~0, \texttt{3{,}3\,V} = bit~1.
Wadą jest podatność na zakłócenia — impuls elektryczny na kablu bezpośrednio
zaburza sygnał.

W~sygnałowaniu \emph{różnicowym} informacja jest zakodowana w~\textbf{różnicy
napięć} między dwoma przewodami: CAN\_H (\textit{CAN High}) i~CAN\_L
(\textit{CAN Low}). Odbiornik mierzy nie napięcie bezwzględne, lecz
$\Delta U = U_{\text{CAN\_H}} - U_{\text{CAN\_L}}$:

\begin{center}
\begin{tabular}{lrrr}
\toprule
\textbf{Stan} & $U_{\text{CAN\_H}}$ & $U_{\text{CAN\_L}}$ & $\Delta U$ \\
\midrule
\textit{recessive} (1) & $\approx 2{,}5$\,V & $\approx 2{,}5$\,V & $\approx 0$\,V \\
\textit{dominant}  (0) & $\approx 3{,}5$\,V & $\approx 1{,}5$\,V & $\approx 2$\,V \\
\bottomrule
\end{tabular}
\end{center}

Zakłócenie elektromagnetyczne (np.\ impuls od silnika elektrycznego) trafia
\emph{jednakowo} na oba przewody biegnące obok siebie, podnosząc napięcie na
obu o~tę samą wartość. Różnica $\Delta U$ pozostaje wtedy niezmieniona.
Stąd nazwa: sygnałowanie \textbf{różnicowe} — bo liczymy \emph{różnicę},
nie wartość bezwzględną.

\bigskip
\textbf{Arbitracja — co to znaczy?}

Arbitraż (z~łac.\ \textit{arbitrare} — rozsądzać) to procedura rozstrzygania
konfliktu o~dostęp do wspólnego zasobu. W~sieci CAN \emph{każdy węzeł może
zacząć nadawać w~dowolnej chwili} — nie ma nadrzędnego kontrolera, który
dzieliłby czas antenowy. Gdy dwa węzły zaczną nadawać jednocześnie, nastąpi
\textbf{kolizja}.

Arbitracja CAN rozwiązuje kolizję \emph{bez niszczenia żadnej ramki}
i~bez~żadnej dodatkowej komunikacji:
\begin{enumerate}[label=\arabic*.]
    \item Każdy nadający węzeł jednocześnie \emph{słucha} magistrali podczas
          własnej transmisji.
    \item Identyfikator nadawany jest od bitu najbardziej znaczącego (MSB).
    \item Jeśli węzeł wysyła bit~1 (\textit{recessive}), ale widzi na magistrali
          bit~0 (\textit{dominant}), oznacza to, że \emph{inny węzeł ma niższy
          numer ID} — czyli wyższy priorytet.
    \item Ten węzeł \textbf{natychmiast przestaje nadawać} i~czeka na kolejną
          okazję.
    \item Węzeł o~najniższym ID nadaje dalej bez zakłóceń — jego ramka
          pozostaje nienaruszona.
\end{enumerate}
Niższy numer ID = wyższy priorytet. Kolizja nie niszczy żadnej wiadomości
(stąd angielskie określenie \textit{non-destructive}).

\bigskip
\textbf{Stan dominant — co to ,,ciągnięcie do 0\,V'' i~,,oba węzły''?}

Magistrala CAN w~stanie \textit{recessive} (spoczynkowym) jest pasywna:
rezystory terminacyjne na obu końcach kabla utrzymują oba przewody na napięciu
$\approx 2{,}5$\,V. Nikt niczego aktywnie nie steruje.

Stan \textit{dominant} wymaga aktywnego działania: nadajnik (układ scalony
zwany \textit{transceiveriem}) włącza wewnętrzne tranzystory, które:
\begin{itemize}
    \item \textbf{ciągną CAN\_L w~dół} — w~kierunku masy (0\,V), obniżając
          napięcie z~$2{,}5$\,V do $\approx 1{,}5$\,V,
    \item \textbf{ciągną CAN\_H w~górę} — w~kierunku napięcia zasilania,
          podnosząc napięcie z~$2{,}5$\,V do $\approx 3{,}5$\,V.
\end{itemize}
Słowo \emph{,,ciągnięcie''} jest standardowym terminem elektronicznym:
tranzystor \emph{ciągnie} przewód do określonego napięcia, ponieważ musi
dostarczyć prąd przezwyciężający rezystory terminacyjne i~pojemności kabla —
podobnie jak ciągnie się linę pod górę wbrew grawitacji.

\textbf{,,Oba węzły ciągną''} odnosi się do topologii szyny.
Wszystkie węzły CAN podłączone są równolegle do tych samych dwóch przewodów.
Kluczowa właściwość: \emph{jeśli choćby jeden węzeł narzuci stan dominant,
cała magistrala przejdzie w~ten stan} — bez względu na to, co robią pozostałe.
Gdy żaden węzeł nie ciągnie (wszystkie pasywne), magistrala sama wraca do
\textit{recessive}. Jest to logika \textbf{iloczynu z~otwartym kolektorem}
(tzw.\ wired-AND): jeden dominant bierze górę nad wszystkimi recessive —
tak jak 0 \& 1 \& 1 = 0.

% ─────────────────────────────────────────────────────────────────────────────
\subsection*{A.1 Sieć CAN — podstawy}

CAN (\textit{Controller Area Network}) to standard komunikacyjny opracowany przez
firmę Bosch w~1986~roku, pierwotnie na potrzeby motoryzacji. Dziś stosowany jest
również powszechnie w~automatyce przemysłowej, medycynie i~robotyce.

\textbf{Kluczowe właściwości:}
\begin{itemize}
    \item \textbf{Magistrala wielowęzłowa} (\textit{multi-master bus}) — każdy
          węzeł może rozpocząć nadawanie w~dowolnej chwili.
    \item \textbf{Sygnałowanie różnicowe} (CAN\_H, CAN\_L) — odporność na
          zakłócenia elektromagnetyczne.
    \item \textbf{Arbitracja bitowa} (\textit{bitwise non-destructive arbitration}) —
          w~przypadku kolizji wygrywa wiadomość o~niższym identyfikatorze
          (wyższym priorytecie), bez uszkodzenia ramki.
    \item \textbf{Dwa stany logiczne:}
          \textit{dominant} (logiczne~0, oba węzły ciągną magistralę do 0~V)
          oraz \textit{recessive} (logiczne~1, magistrala w~stanie biernym).
          Przy kolizji stan dominant zawsze wygrywa.
    \item Prędkość: do 1\,Mb/s (klasyczny CAN), do 8\,Mb/s (CAN~FD).
\end{itemize}

\textbf{Priorytety:} identyfikator wiadomości pełni równocześnie rolę priorytetu.
Węzeł o~niższym numerze~ID (więcej bitów dominant na początku) ma wyższy
priorytet. Dzięki temu w~sieci CAN wiadomości krytyczne czasowo (np.\ dane z~czujników
bezpieczeństwa) zawsze dostają się na magistralę przed wiadomościami o~niższym priorytecie.

% ─────────────────────────────────────────────────────────────────────────────
\subsection*{A.2 Struktura rozszerzonej ramki CAN (CAN~2.0B)}

Specyfikacja CAN~2.0 definiuje dwa formaty ramek danych:
\begin{itemize}
    \item \textbf{Standardowy} (CAN~2.0A) — 11-bitowy identyfikator,
    \item \textbf{Rozszerzony} (CAN~2.0B) — 29-bitowy identyfikator.
\end{itemize}

Zadanie 4.3 dotyczy ramki rozszerzonej z~8~bajtami danych.
Poniższa tabela opisuje każde pole ramki:

\begin{table}[H]
\centering
\small
\begin{tabular}{llrp{6.5cm}}
\toprule
\textbf{Pole} & \textbf{Podpole} & \textbf{Bitów} & \textbf{Opis} \\
\midrule
SOF & & 1 & \textit{Start of Frame} — bit dominant otwierający ramkę \\
\midrule
\multirow{5}{*}{Arbitration} & Base ID [28:18] & 11 & Starsze 11~bitów 29-bitowego ID \\
 & SRR & 1 & \textit{Substitute Remote Request} — zawsze recessive \\
 & IDE & 1 & \textit{Identifier Extension} — recessive = ramka rozszerzona \\
 & Extension ID [17:0] & 18 & Młodsze 18~bitów 29-bitowego ID \\
 & RTR & 1 & \textit{Remote Transmission Request} \\
\midrule
\multirow{3}{*}{Control} & r1 & 1 & Bit zarezerwowany (dominant) \\
 & r0 & 1 & Bit zarezerwowany (dominant) \\
 & DLC [3:0] & 4 & \textit{Data Length Code} (0\,–\,8) \\
\midrule
Data & & 64 & 8~bajtów danych użytecznych \\
\midrule
\multirow{2}{*}{CRC} & CRC sequence & 15 & Wartość sumy kontrolnej CRC-15 \\
 & CRC delimiter & 1 & Bit recessive kończący pole CRC \\
\midrule
\multirow{2}{*}{ACK} & ACK slot & 1 & Odbiornik wstawia bit dominant \\
 & ACK delimiter & 1 & Bit recessive \\
\midrule
EOF & & 7 & \textit{End of Frame} — 7~bitów recessive \\
\midrule
IFS & & 3 & \textit{Interframe Space} — 3~bity recessive \\
\bottomrule
\end{tabular}
\caption{Pola rozszerzonej ramki danych CAN~2.0B przy 8~bajtach danych.}
\end{table}

\textbf{Podsumowanie bitowe} dla $\text{DLC} = 8$ (8~bajtów danych):

\begin{align*}
L_{\text{stały}} &= 1 + (11+1+1+18+1) + (1+1+4) + 64 + (15+1) + (1+1) + 7 + 3 \\
                    &= 1 + 32 + 6 + 64 + 16 + 2 + 7 + 3 = \mathbf{131\ \text{bitów}}
\end{align*}

\textbf{Region bit stuffingu} (SOF aż do końca sekwencji CRC, bez delimitera):

\begin{equation}
N_{\text{fill}} = 1 + 32 + 6 + 64 + 15 = \mathbf{118\ \text{bitów}}
\end{equation}

% ─────────────────────────────────────────────────────────────────────────────
\subsection*{A.3 Technika bit stuffing — zasada i~zastosowanie}

\textbf{Zasada:} po każdych 5~kolejnych bitach \emph{tej samej} polaryzacji (same dominy
lub same recessive) nadajnik automatycznie wstawia jeden bit o~\emph{przeciwnej}
polaryzacji (bit stuffujący, ang.\ \textit{stuff bit}). Odbiornik usuwa te dodatkowe
bity ze strumienia danych.

\textbf{Cel:} zapewnienie wystarczającej liczby przejść sygnału w~strumieniu bitowym,
co umożliwia odbiornikowi synchronizację zegara (\textit{clock recovery}) bez
osobnego sygnału taktowania.

\textbf{Zakresy stosowania:}
\begin{itemize}
    \item Bit stuffing \textbf{obowiązuje}: SOF, pole arbitracji, pole sterowania,
          pole danych, sekwencja CRC-15.
    \item Bit stuffing \textbf{nie obowiązuje}: delimiter CRC, pole ACK, EOF, IFS
          (te pola mają sztywno określoną wartość bitu).
\end{itemize}

\textbf{Przykład:}
\begin{center}
\begin{tabular}{ll}
\toprule
Strumień wejściowy  & \texttt{1 1 1 1 1 0 0 0 0 0 1 1} \\
Strumień po stuffingu & \texttt{1 1 1 1 1 \textbf{[0]} 0 0 0 0 0 \textbf{[1]} 1 1} \\
\bottomrule
\end{tabular}
\end{center}
Bity stuffujące zaznaczono nawiasami $[\,]$.

\textbf{Maksymalna liczba bitów stuffujących} dla regionu o~$N$ bitach:
\begin{equation}
N_{\text{stuff,max}} = \left\lfloor \frac{N_{\text{fill}}}{5} \right\rfloor
\end{equation}
Dla $N_{\text{fill}} = 118$:
\begin{equation}
N_{\text{stuff,max}} = \left\lfloor \frac{118}{5} \right\rfloor = \mathbf{23}
\end{equation}

\textbf{Minimalna liczba bitów stuffujących} wynosi 0 — osiągalna, gdy strumień
nie zawiera 5~kolejnych bitów tej samej polaryzacji (np.\ wzorzec naprzemienny
\texttt{10101010\ldots}).

% ─────────────────────────────────────────────────────────────────────────────
\subsection*{A.4 Identyfikatory CAN i~ograniczenia specyfikacji}

\textbf{Teoretyczna przestrzeń identyfikatorów:}
\begin{itemize}
    \item Ramka standardowa (11-bitowa): $2^{11} = 2048$
    \item Ramka rozszerzona (29-bitowa): $2^{29} = 536\,870\,912$
\end{itemize}

\textbf{Ograniczenie specyfikacji:} identyfikatory, których 7~najbardziej znaczących
bitów jest \emph{jednocześnie} w~stanie recessive (wszystkie~1), są zarezerwowane
i~nie mogą być używane. Powodem jest możliwość pomyłki takiej ramki
z~\textit{passive error flag} (6~bitów recessive) lub polem IFS.

\begin{table}[H]
\centering
\begin{tabular}{lrrr}
\toprule
\textbf{Format} & \textbf{Bity ID} & \textbf{Zakazane ID} & \textbf{Dostępne ID} \\
\midrule
Standardowy & 11 & $2^{4} = 16$ \quad (ID[10:4] = 1111111) & $2048 - 16 = \mathbf{2032}$ \\
Rozszerzony & 29 & $2^{22} = 4\,194\,304$ \quad (ID[28:22] = 1111111) & $536\,870\,912 - 4\,194\,304 = \mathbf{532\,676\,608}$ \\
\bottomrule
\end{tabular}
\caption{Rzeczywista liczba dostępnych identyfikatorów w~sieci CAN.}
\end{table}

\textbf{Wniosek:} w~praktyce dostępna przestrzeń identyfikatorów jest nieznacznie
mniejsza niż teoretyczna. W~typowych sieciach przemysłowych ($\leq$100~węzłów)
ograniczenie to nie ma praktycznego znaczenia.

% ─────────────────────────────────────────────────────────────────────────────
\subsection*{A.5 Suma kontrolna CRC-15/CAN — teoria i~gwarancje}

\subsubsection*{Idea CRC}

CRC (\textit{Cyclic Redundancy Check}) to suma kontrolna wyznaczana przez traktowanie
ciągu bitów jako współczynników wielomianu nad ciałem $\text{GF}(2)$
(dodawanie = XOR, mnożenie = AND, brak przeniesień).

Nadajnik dzieli wielomian wiadomości $M(x)$ przez wybrany wielomian generujący $G(x)$
stopnia~15 i~dołącza resztę $R(x)$ do transmitowanej ramki:
\[
R(x) = M(x) \cdot x^{15} \bmod G(x)
\]
Odbiornik sprawdza, czy odebrany strumień jest podzielny przez $G(x)$ bez reszty.
Każda niedająca się podzielić reszta wskazuje na błąd transmisji.

\subsubsection*{Wielomian generujący CRC-15/CAN}

\begin{equation}
G(x) = x^{15} + x^{14} + x^{10} + x^8 + x^7 + x^4 + x^3 + 1
\end{equation}

Reprezentacja szesnastkowa (bez bitu $x^{15}$, używana w~algorytmie):
\[
G_{\text{hex}} = \texttt{0x4599}
= \underbrace{0100\ 0101\ 1001\ 1001}_{\text{bity 14\ldots0}}
\]

\subsubsection*{Gwarancje wykrywania błędów}

Wielomian CRC-15/CAN został starannie dobrany przez Boscha, aby dla typowych
rozmiarów ramek CAN gwarantować:
\begin{itemize}
    \item wykrycie \textbf{wszystkich błędów 1-bitowych},
    \item wykrycie \textbf{wszystkich błędów 2-bitowych} (dla dowolnej długości ramki),
    \item wykrycie \textbf{wszystkich serii błędów} o~długości $\leq 15$~bitów,
    \item wykrycie \textbf{wszystkich nieparzystych liczb błędów},
    \item prawdopodobieństwo niewykrycia błędu nieobjętego powyższymi gwarancjami:
          $\approx 3{,}1 \times 10^{-5}$ dla typowej długości ramki CAN.
\end{itemize}

\subsubsection*{Algorytm LFSR z~zadania}

Algorytm opisany w~treści zadania to implementacja LFSR
(\textit{Linear Feedback Shift Register} — rejestr przesuwający z~liniowym sprzężeniem
zwrotnym). Przetwarza bity od SOF do ostatniego bitu pola danych (włącznie),
\textbf{bit po bicie, począwszy od MSB}:

\begin{center}
\begin{tabular}{ll}
\toprule
\textbf{Zmienna} & \textbf{Znaczenie} \\
\midrule
\texttt{CRC\_RG(14:0)} & 15-bitowy rejestr CRC (bit 14 = MSB) \\
\texttt{NXTBIT} & kolejny bit wejściowy (0 lub 1) \\
\texttt{CRCNXT} & bit pomocniczy = XOR bitu wejściowego z~MSB rejestru \\
\bottomrule
\end{tabular}
\end{center}

\begin{enumerate}
    \item \texttt{CRC\_RG = 0} (inicjalizacja)
    \item Dla każdego bitu wejściowego \texttt{NXTBIT}:
    \begin{enumerate}[label=\alph*)]
        \item $\texttt{CRCNXT} = \texttt{NXTBIT} \oplus \texttt{CRC\_RG}(14)$
        \item przesuń rejestr w~lewo o~1: $\texttt{CRC\_RG}(14:1) \leftarrow \texttt{CRC\_RG}(13:0)$;
              $\texttt{CRC\_RG}(0) = 0$
        \item jeśli $\texttt{CRCNXT} = 1$: $\texttt{CRC\_RG}(14:0) \mathrel{\oplus}= \texttt{0x4599}$
    \end{enumerate}
    \item Wynik to zawartość \texttt{CRC\_RG} po przetworzeniu wszystkich bitów.
\end{enumerate}

Jest to równoważne dzieleniu wielomianu wiadomości przez $G(x)$ w~ciele $\text{GF}(2)$.


% ━━━━━━━━━━━━━━━━━━━━━━━━━━━━━━━━━━━━━━━━━━━━━━━━━━━━━━━━━━━━━━━━━━━━━━━━━━━━
\section*{Dodatek B — Wymagana wiedza i~kluczowe aspekty przy implementacji}
% ━━━━━━━━━━━━━━━━━━━━━━━━━━━━━━━━━━━━━━━━━━━━━━━━━━━━━━━━━━━━━━━━━━━━━━━━━━━━

% ─────────────────────────────────────────────────────────────────────────────
\subsection*{B.1 Arytmetyka nad GF(2) — dlaczego XOR jest sercem CRC}

Ciało $\text{GF}(2) = \{0, 1\}$ ma następujące działania:
\[
0 \oplus 0 = 0,\quad 1 \oplus 0 = 1,\quad 1 \oplus 1 = 0
\quad\Longrightarrow\quad
\text{dodawanie = XOR, odejmowanie = XOR, brak przeniesień}
\]

Mnożenie: $0 \cdot 0 = 0$, $1 \cdot 0 = 0$, $1 \cdot 1 = 1$ (AND).

Ciąg bitów $b_{n-1}b_{n-2}\ldots b_1 b_0$ reprezentuje wielomian:
$b_{n-1}x^{n-1} + b_{n-2}x^{n-2} + \ldots + b_0$.

Przesunięcie rejestru w~lewo o~1 = mnożenie przez $x$.
XOR z~wielomianem generującym = odjęcie wielomianu (w~GF(2) dodawanie = odejmowanie).

Dlatego algorytm LFSR (przesuń + warunkowo XOR z~generatorem) to \emph{dokładnie}
kolejne kroki dzielenia wielomianu przez $G(x)$ w~GF(2).

% ─────────────────────────────────────────────────────────────────────────────
\subsection*{B.2 Rejestr LFSR — jak algorytm z~zadania działa w~praktyce}

Algorytm opisany w~treści zadania jest implementacją \textbf{LFSR}
(\textit{Linear Feedback Shift Register} — liniowy rejestr przesuwający
z~sprzężeniem zwrotnym). Można go wyobrazić sobie jako 15-bitowy rejestr
sprzętowy, w~którym każdy takt zegara:
\begin{enumerate}[label=\arabic*.]
    \item przesuwa wszystkie bity o~jedną pozycję w~lewo (bit~14 wypada),
    \item na miejsce bitu~0 wstawia~0,
    \item jeśli wypadły bit~14 różni się XOR-em od bitu wejściowego,
          dokonuje XOR całego rejestru z~wielomianem \texttt{0x4599}.
\end{enumerate}

Kluczowe typy danych w~C++: rejestr to \texttt{uint16\_t} (16~bitów bez znaku),
przy czym używamy tylko dolnych 15~bitów — po każdym przesunięciu maskujemy
wartość przez \texttt{\& 0x7FFF}. Wejście wczytujemy jako ciąg znaków \texttt{'0'}/\texttt{'1'};
znaki ASCII \texttt{'0'} i~\texttt{'1'} różnią się tylko bitem~0, więc operacja
\texttt{znak \& 1} daje bezpośrednio wartość bitu bez żadnej konwersji.

% ─────────────────────────────────────────────────────────────────────────────
\subsection*{B.3 Strategie optymalizacji algorytmu}

\subsubsection*{B.3.1 Eliminacja instrukcji warunkowej (\textit{branchless})}

Wewnętrzna pętla algorytmu zawiera instrukcję \texttt{if}: jeśli \texttt{CRCNXT}~=~1,
wykonaj XOR z~wielomianem; w~przeciwnym razie pomiń. Przy losowych danych wejściowych
procesor myli się w~przewidywaniu gałęzi (\textit{branch misprediction}) mniej więcej
co drugi takt, tracąc 5\,–\,15~cykli na każdą pomyłkę.

Rozwiązanie: zastąpić \texttt{if} operacją arytmetyczną lub maską bitową, tak aby
wynik XOR był zero gdy \texttt{CRCNXT}~=~0 i~równy wielomianowi gdy \texttt{CRCNXT}~=~1.
Mnożenie \texttt{CRCNXT * wielomian} oraz maska z~negacji jednoargumentowej
(\texttt{-CRCNXT} daje \texttt{0x0000} lub \texttt{0xFFFF}) to dwie typowe techniki.
Przy \texttt{-O3} kompilator często sam wykrywa ten wzorzec i~generuje bezgałęziowy kod.

\subsubsection*{B.3.2 Tablica przeglądowa (\textit{lookup table})}

Zamiast przetwarzać wejście bit po bicie (96~iteracji dla 96~bitów), można
przetworzyć je \textbf{bajt po bajcie} (12~iteracji dla 96~bitów = 8-krotne
przyspieszenie iteracyjne). Wymaga to prekomputacji 256-elementowej tablicy,
w~której \texttt{TABLE[i]} przechowuje stan rejestru po przetworzeniu bajtu
o~wartości~\texttt{i} startując od pustego rejestru.

Dla CRC-15 (MSB-first, rejestr 15-bitowy) klucz do tablicy to XOR górnych
8~bitów aktualnego rejestru (bity~14\ldots7) z~wejściowym bajtem. Tablica ma
rozmiar 256\,×\,2\,B = 512\,B, co mieści się w~całości w~pamięci podręcznej L1 (32\,–\,64\,kB),
eliminując kosztowne odczyty z~RAM po rozgrzaniu pętli.

Dzięki słowu kluczowemu \texttt{constexpr} tablica może być wygenerowana
\emph{w~czasie kompilacji} i~umieszczona bezpośrednio w~sekcji \texttt{.rodata}
pliku wykonywalnego — zero kosztu inicjalizacji w~czasie uruchomienia.

\subsubsection*{B.3.3 Flagi kompilatora}

Kompilacja z~\texttt{-O3 -march=native} włącza agresywną optymalizację:
rozwijanie pętli (\textit{loop unrolling}), \textit{inlining} funkcji,
reordering instrukcji pod potok (\textit{pipeline}) procesora oraz instrukcje
wektorowe specyficzne dla danego CPU (NEON na Apple Silicon, AVX2 na x86-64).
Dla krótkiej pętli CRC (96~iteracji) kompilator może całkowicie rozwinąć pętlę
do ciągu instrukcji bez skoków.

% ─────────────────────────────────────────────────────────────────────────────
\subsection*{B.4 Pomiar czasu i~integralność benchmarku}

\subsubsection*{B.4.1 Mierzenie czasu}

Do pomiaru czasu służy \texttt{std::chrono::high\_resolution\_clock}
z~nagłówka \texttt{<chrono>}. Zegar ten ma rozdzielczość rzędu nanosekund
na macOS i~Linux. Wzorzec użycia: pobrać punkt startowy przed pętlą,
punkt końcowy po pętli, a~różnicę rzutować na żądaną jednostkę
(\texttt{milliseconds}, \texttt{nanoseconds} itp.) metodą \texttt{duration\_cast}.
Dla wiarygodnych wyników $n$ powinno wynosić co najmniej $10^6$ — przy mniejszych
wartościach dominuje szum pomiaru zegara systemowego.

\subsubsection*{B.4.2 Pułapka optymalizatora — eliminacja martwego kodu}

Przy \texttt{-O3} kompilator może wykryć, że wynik pętli benchmarkowej
($n$~wywołań tej samej funkcji z~tym samym argumentem) jest stały i~zastąpić
całą pętlę \emph{jednym} wywołaniem. Benchmark pokazałby wówczas czas
jednej iteracji niezależnie od wartości~$n$.

Dwa mechanizmy zapobiegające temu zachowaniu:
\begin{itemize}
    \item \textbf{\texttt{volatile} sink} — zapis wyniku do zmiennej
          \texttt{volatile} jest przez standard C++ traktowany jako
          \emph{obserwowalny efekt uboczny}; kompilator nie może go usunąć
          ani wynieść poza pętlę (\textit{loop-invariant code motion}).
    \item \textbf{\texttt{\_\_attribute\_\_((noinline))} / \texttt{\_\_declspec(noinline)}}
          na funkcji CRC — uniemożliwia kompilatorowi wgląd w~ciało funkcji
          przez granicę translacji, co blokuje wszelką analizę inwariantów pętli.
\end{itemize}

Oba mechanizmy należy stosować jednocześnie jako niezależne zabezpieczenia
integralności benchmarku — analogicznie jak w~projekcie \texttt{modbus-rtu-crc}.

% ─────────────────────────────────────────────────────────────────────────────
\subsection*{B.5 Wyznaczanie $T_{\min}$ i~$T_{\max}$ — metoda obliczeniowa (zadanie 4.3)}

\subsubsection*{Parametry bazowe}

Przy prędkości 1\,Mb/s czas jednego bitu wynosi $T_b = 1\,\mu$s.

Liczba \emph{stałych} bitów ramki (bez bitów stuffujących), patrz tabela w~A.2:

\begin{align*}
L_{\text{stały}} &= 1 + 32 + 6 + 64 + 16 + 2 + 7 + 3 = 131\ \text{bitów}
\end{align*}

Rozmiar regionu stuffowanego:
\begin{align*}
N_{\text{fill}} &= 1 + 32 + 6 + 64 + 15 = 118\ \text{bitów}
\end{align*}

\subsubsection*{Najkrótszy czas — minimum bitów stuffujących (0)}

Najkrótsza ramka nie zawiera żadnych bitów stuffujących. Jest to możliwe,
gdy 118-bitowy region stuffowany nie posiada żadnej sekwencji 5~kolejnych bitów
tej samej polaryzacji.

\textbf{Przykładowe wejście osiągające 0~bitów stuffujących:}
\begin{itemize}
    \item Dane (8~bajtów): \texttt{0x55 0x55 0x55 0x55 0x55 0x55 0x55 0x55}
          (\texttt{0x55} = \texttt{01010101} — wzorzec naprzemienny)
    \item Identyfikator 29-bitowy: np.\ \texttt{0x0AAAAAA}
          (naprzemienne bity w~ID)
\end{itemize}

\begin{align*}
L_{\min} &= L_{\text{stały}} + 0 = \mathbf{131\ \text{bitów}} \\
T_{\min} &= 131 \cdot T_b = \mathbf{131\ \mu\text{s}}
\end{align*}

Współczynnik efektywności dla najkrótszej ramki:
\begin{equation}
\eta_{\max} = \frac{64}{131} \approx \mathbf{48{,}9\%}
\end{equation}

\subsubsection*{Najdłuższy czas — maksimum bitów stuffujących (23)}

Najdłuższa ramka ma maksymalną liczbę bitów stuffujących.
Warunek: co 5~bitów w~regionie stuffowanym musi być tego samego stanu.

Wzorzec generujący maksi\-malne stuffowanie:
\[
\underbrace{11111}_{5} \underbrace{00000}_{5} \underbrace{11111}_{5} \underbrace{00000}_{5} \cdots
\]
Dla 118~bitów: $\lfloor 118 / 5 \rfloor = 23$ bity stuffujące.

\textbf{Przykładowe dane (dane pola 8~bajtów):}
Wartość \texttt{0x00 0x00 0x00 0x00 0x00 0x00 0x00 0x00} (same zera w~polu danych)
daje 64~bity~\texttt{0} z~rzędu $\Rightarrow$ $\lfloor 64/5 \rfloor = 12$~bitów stuffujących
tylko z~pola danych, a~odpowiednio dobrany~ID + równanie CRC mogą uzupełnić do 23.

\textbf{Uwaga praktyczna:} znalezienie dokładnego zestawu ID + danych prowadzącego
do \emph{dokładnie} 23~bitów stuffujących wymaga uwzględnienia wyniku CRC
(który jest wyznaczany przez ID + dane i~nie jest swobodnie wybierany). Konieczne
jest iteracyjne sprawdzanie kandydatów.

\begin{align*}
L_{\max} &= L_{\text{stały}} + 23 = 131 + 23 = \mathbf{154\ \text{bitów}} \\
T_{\max} &= 154 \cdot T_b = \mathbf{154\ \mu\text{s}}
\end{align*}

Współczynnik efektywności dla najdłuższej ramki:
\begin{equation}
\eta_{\min} = \frac{64}{154} \approx \mathbf{41{,}6\%}
\end{equation}

\subsubsection*{Podsumowanie zadania 4.3}

\begin{center}
\begin{tabular}{lrrrr}
\toprule
\textbf{Scenariusz} & \textbf{Bity stałe} & \textbf{Bity stuffuj.} & \textbf{Bity całk.} & $\bm{\eta}$ \\
\midrule
Najkrótsza ramka & 131 & 0  & 131 & 48,9\% \\
Najdłuższa ramka & 131 & 23 & 154 & 41,6\% \\
\bottomrule
\end{tabular}
\end{center}

Przy 1\,Mb/s: $T_{\min} = 131\,\mu$s, $T_{\max} = 154\,\mu$s.

\end{document}
